% Added 27 November 2026 for Edward B. Baker III. Drafted by Claude Fable 5.2
% (Anthropic, Cowork session) from the 08 tube investigation's notes E39, E70,
% E76 and their replays E73 and E79 at b607bd1, with an independent replay of
% the group computation filed with the brainstorm note
% ALTERNATIVE_FRAMEWORKS_AND_INVESTIGATIONS_2_20261127.md (check_05).
\section{A recognized fibred closing: the even prism sphere}\label{sec:prism}

The tube investigation of the cut-geometry program produced one smooth
recognition that does not use curvature. We record its statement, the
structure of its proof and its present level of verification, because it
is an instance of the problem this paper is about: a homotopy four-sphere
defined by surgery whose smooth type was not supplied by the sources
consulted.

Write \(\mathbb H\) for the quaternions and \(S^3\) for the unit
quaternions. Let \(z=e^{2\pi i/17}\) and let
\(\Gamma=Q_8\times C_{17}\) act freely on \(S^3\) by
\((u,z^r)\cdot q=u\,q\,z^{-r}\), with quotient the prism manifold
\(Y=S^3/\Gamma\). Put \(b=\tfrac12(1+i+j+k)\). The isometry
\(q\mapsto b\,q\,j^{-1}\) normalizes \(\Gamma\), since conjugation by
\(b\) cycles \(i,j,k\) and conjugation by \(j\) inverts \(z\). It induces
an isometry \(J\) of \(Y\) of order six. Comparing real parts of
\(u^{-1}b^m\) and \(z^{-r}j^m\) shows that \(J\) and \(J^2\) act freely
on \(Y\), while \(J^3\) fixes circles. Let
\[
 N=Y\times[0,1]\big/\,(y,1)\sim(Jy,0)
\]
be the mapping torus. Choose \(p_0\in S^3\) with
\(p_0\,j\,p_0^{-1}=(i+j+k)/\sqrt3\). Then
\(b\,p_0=p_0e^{j\pi/3}\), so \(J(p_0e^{j\pi/6})=p_0\) and
\[
 \gamma(t)=\bigl[\,p_0e^{j\pi t/6},\,t\,\bigr],\qquad 0\le t\le1,
\]
is a closed section of \(N\to S^1\). The translation flow of \(N\) is a
smooth circle action without fixed points. Its orbit space is the
spherical orbifold \(O=Y/\langle J\rangle\), whose underlying space is
\(L(3,1)\) and whose singular set is one circle of order two. The image
\(\ell\) of \(\gamma\) is a closed geodesic in the regular part of
\(O\), and it is a core of a genus-one Heegaard splitting of
\(L(3,1)\). Let \(\mu\) be its meridian. Filling the exterior of
\(\ell\) in the underlying space along a longitude gives a lens space,
and exactly one longitude \(a\) gives \(S^3\). A longitude of \(\ell\)
together with the orbit direction frames the normal bundle of
\(\gamma\). The longitudes \(a+s\mu\) with \(s\) even give one of the two
framings of \(\gamma\), which we call even, and those with \(s\) odd give
the other.

\begin{proposition}[The even prism closing]\label{prop:even-prism}
Let \(\Sigma_{\rm ev}\) and \(\Sigma_{\rm odd}\) be the results of
surgery on \(\gamma\subset N\) in the even and the odd framing.
Both are smooth homotopy four-spheres. The even one is diffeomorphic to
\(S^4\). Consequently the belt sphere \(K\subset\Sigma_{\rm ev}=S^4\)
is a smooth fibred two-knot with fibre the punctured prism manifold
\(Y\setminus{\rm int}\,B^3\), monodromy of order six on the closed fibre
and group \(\Gamma\rtimes\mathbb Z\), and \(\Sigma_{\rm odd}\) is the
Gluck twist of \(S^4\) along \(K\).
\end{proposition}

\begin{proof}[Structure of the proof]
The group of \(N\) is \(\Gamma\rtimes\mathbb Z\) with the class of
\(\gamma\) as stable letter. Killing it identifies \(i,j,k\), hence
kills \(Q_8\), and identifies \(z\) with \(z^{-1}\), hence kills
\(C_{17}\). Each surgery is therefore simply connected with Euler
characteristic two, which is the realization of Plotnick and Suciu
\cite[Proposition 5.1]{PlotnickSuciu1987}. The two attaching maps
differ by the nontrivial loop of \(SO(3)\), which gives the last
assertion once the even closing is identified.

The preimage of a tube about \(\ell\) is
\(S^1_a\times D^2_\mu\times S^1_f\), where \(f\) is the orbit direction,
and the orbit of \(\gamma\) is the torus \(T=S^1_a\times\{0\}\times S^1_f\).
In \(\Sigma_{\rm ev}\) the surgery compresses \(T\) along the disk
bounded by \(a\) to an embedded two-sphere \(S_T\) with a product
neighbourhood. Surgery on \(S_T\) replaces the preimage of the tube by
\(S^1_\mu\times D^2_a\times S^1_f\). The result \(W\) carries a
fixed-point-free circle action extending the one on the exterior, and
its orbit orbifold is the filling \(O(a)\), in which the new core is
regular and the circle of order two is unchanged. Conversely
\(\Sigma_{\rm ev}\) is recovered from \(W\) by a framed surgery on the
core circle \(\delta\) of the new piece.

The computed input is the orbifold group
\(\pi_1^{\rm orb}O(a)\cong S_3\). It is obtained from a
four-generator presentation of the orbifold group of the exterior of
\(\ell\) read from a tangle diagram
\cite[TUBE, equivariant surgery]{Baker2026MarkedCutCalculations}.
The Thom isomorphism for \(S_T\) gives \(H_1(W)=\mathbb Z\), generated
by \(\delta\). The principal orbit \(h\) is central in
\(\Pi=\pi_1(W)\) and \(\Pi/\langle h\rangle\cong S_3\). Abelianizing
shows that \(h\) has degree \(\pm2\) and infinite order, so the
commutator subgroup of \(\Pi\) is \(C_3\) and
\(\Pi\cong C_3\rtimes\mathbb Z\) with a degree-one element acting by
inversion. Since squaring is a bijection of \(C_3\), all elements of
degree one are conjugate. The exceptional orbit \(e\) of order two,
suitably oriented, has degree one. Hence \(\delta\) is freely homotopic,
and in dimension four isotopic, to \(e\).

Near \(e\) the action is
\(u\cdot(w,y)=(u^{m}w,R_{k\arg u}\,y)\) on \(S^1\times D^3\), with
\(|m|=2\), \(k\) odd and \(R\) the rotation about the slice axis.
Rotating the product frame once changes the framing class and replaces
\(k\) by \(k+m\). In either class the same formula on \(D^2\times S^2\)
extends the action across the surgery. Therefore \(\Sigma_{\rm ev}\)
carries a smooth effective circle action. A closed simply connected
smooth four-manifold with a circle action is a connected sum of copies
of \(S^4\), \(\pm\mathbb{CP}^2\) and \(S^2\times S^2\)
\cite{Fintushel1976Circle,Pao1978Nonlinear},
\cite[Theorem 3.2]{GalazGarciaSearle2014}, and \(b_2=0\) leaves \(S^4\).
\end{proof}

\paragraph{Verification status.}
The argument is a working proof in the investigation's notes, with two
independent model replays of the group and of the transfer
\cite[TUBE, marked surgery and its independent replay]{Baker2026MarkedCutCalculations}.
For this section the presentation was replayed by a separate
coset enumeration. With the meridian of \(\ell\) killed it gives a group
of order \(816\), and an explicit surjection identifies it with the
orbifold group \((2T\times{\rm Dic}_{17})/\{\pm1\}\) of \(O\). At the
longitude \(a\) it gives a group of order six. With the two meridians of
order two also killed, the filling along \(a+s\mu\) has order
\(|3s-1|\), in agreement with the lens spaces above. The tangle diagram
from which the presentation is read was not rederived, the proposition
has not been examined by a topologist, and it has no formal proof. Its
formalization needs the presentation with its diagram certificate, the
group lemmas, and the classification of circle actions as a named
imported theorem.

\paragraph{Relation to earlier work, and what remains.}
Plotnick and Suciu left open whether these homotopy spheres are
standard. For the fibres \(S^3/(Q_8\times C_d)\), fibred two-knots in
\(S^4\) were obtained for \(d=5,11,13,19\) by twisting symmetry-spins of
two-bridge knots \cite{Teragaito1989Twisting}, and the list reported in
\cite{Jablonowski2026Prism} is \(d=3,5,11,13,19,21,27\).
Jab\l{}onowski's recent preprint treats \(d=7\) and states that both
closings are standard, by a scheme of the same kind: the orbit orbifold
of the circle action, fillings whose orbifold groups are \(S_3\) and
\(\mathbb Z/8\), and the classification of circle actions
\cite{Jablonowski2026Prism}. Proposition~\ref{prop:even-prism} applies
that scheme to \(d=17\), which appears in neither list. Its second half
does not transfer. At the adjacent longitudes \(a+\mu\) and \(a-\mu\)
the orbifold groups have abelianizations \(C_4\) and \(C_8\). The
investigation's computations find both groups infinite, and the one
further curve tested there fails the same criterion. The smooth type of
\(\Sigma_{\rm odd}\) is therefore open here. It is an instance of the
Gluck twist problem. Pao's theorem settles that problem for branched
twist spins \cite{Pao1978Nonlinear,Fukuda2018Gluck}, and the closed
monodromy of \(K\) has no fixed point. Whether the group \(S_3\) at the
longitude \(a\) persists for the other odd orders prime to three was not
examined. The proposition is topological. It supplies no positive
subcritical class on a new smooth type and does not bear on entry.
